Un sélecteur de vitesse est formé de deux plaques parallèles séparées de \qty{\v{d}}{mm} entre lesquelles on applique une différence de potentiel de \qty{\v{dV}}{V}. Un champ magnétique uniforme de \qty{\v{B}}{mT}, perpendiculaire au champ électrique et à la vitesse des particules, règne dans la même région. Un faisceau de protons ($m_p = \qty{1.673e-27}{kg}$, $e = \qty{1.602e-19}{C}$) y entre perpendiculairement aux deux champs.

\begin{center}
\begin{tikzpicture}[thick]
    \draw[very thick] (0,1.2) -- (5,1.2);
    \draw[very thick] (0,-1.2) -- (5,-1.2);
    \foreach \x in {0.5,1.5,...,4.5} {
        \node[above] at (\x,1.2) {\small$+$};
        \node[below] at (\x,-1.2) {\small$-$};
    }
    \foreach \x in {0.6,1.3,...,4.8}
        \foreach \y in {-0.8,-0.35,0.35,0.8}
            \node[gray] at (\x,\y) {\small$\times$};
    \node at (2.5,-1.95) {\small$\vec B$ entrant (vers la page)};
    \fill (-1.2,0) circle (2pt) node[below] {\small$p^+$};
    \draw[-{Stealth[length=3mm]}] (-1.2,0) -- node[above] {$\vec v$} (-0.2,0);
    \draw[thin,dashed] (-0.2,0) -- (5.6,0);
    \draw[thin,|-|] (5.9,-1.2) -- node[right] {\small\qty{\v{d}}{mm}} (5.9,1.2);
\end{tikzpicture}
\end{center}

\begin{enumerate}[label=\alph*)]
    \item Déterminez la grandeur du champ électrique entre les plaques.
    \item Déterminez la vitesse des protons qui traversent le sélecteur sans être déviés.
    \item À la sortie du sélecteur, ces protons entrent dans une région où règne seulement le même champ magnétique. Déterminez le rayon de leur trajectoire.
    \item Un proton un peu plus rapide entre dans le sélecteur. De quel côté sera-t-il dévié : du côté où la force électrique le pousse ou du côté où la force magnétique le pousse ? Justifiez.
\end{enumerate}
